\uuid{zJyd}
\titre{Étude de fonction}
\auteur{bodin ; hanani ; mzari}
\organisation{exo7}
\licence{CC BY-NC-SA 4.0}
\datecreate{2020-01-20}
\niveau{L1}
\module{Analyse}
\chapitre{Fonctions circulaires et hyperboliques inverses}
\sousChapitre{Autre}
\difficulte{2}
\theme{Fonctions usuelles, étude de fonction}
\exo7qcmid{lille-1.fonctions-usuelles.018}

\contenu{
\begin{question}
\qclassification{126.00}
\qtags{motcle=étude de fonction}
Soit $f(x)= (-1)^{E(x)}$, où  $E(x)$ est la partie entière de $x$. Quelles sont les assertions vraies ?
\begin{answers}
    \bad{$f$ est périodique de période $1$.}

    \good{$f$ est périodique de période $2$.}

    \bad{$f$ est une fonction paire.}

    \good{$f$ est bornée.}
\end{answers}
\begin{explanations}
$|f(x)|=1$, donc $f$ est bornée.

$\forall x \in \Rr, n\in \Zz, \, E(x+n)= E(x)+n$. On déduit que  $f(x+1)=-f(x)$ et $f(x+2)=f(x)$, donc $f$ est de période $2$.

On a : $ \forall x \in \Rr, \, E(x)\le x< E(x)+1$,  donc $\forall x \in \Rr, \, -E(x)-1< -x\le - E(x)$, on déduit que : $E(-x)=\left\{\begin{array}{cc}-x,& \mbox{si} \, \, x \in \Zz \\ -E(x)-1,& \mbox{si} \,  x \notin \Zz  \end{array}\right.$. Il en découle :  $f(-x)=\left\{\begin{array}{cc}f(x),& \mbox{si} \, \, x \in \Zz \\ -f(x),& \mbox{si} \,  x \notin \Zz  \end{array}\right.$. Donc $f$ n'est ni paire ni impaire.
\end{explanations}
\end{question}
}
